% !TeX root = part9.tex
\documentclass[../final.tex]{subfiles}

\begin{document}

\practice{9}{29.10.2026}{Следы \texorpdfstring{$\gamma$}{γ}-матриц. Спиноры Дирака}

\rtheory{След произведения \texorpdfstring{$\gamma$}{γ}-матриц}

След произведения $n$ матриц сводится к следам произведений $n-2$ матриц:
\begin{equation*}
    \operatorname{Tr}\left[\gamma^{\mu_1}\ldots\gamma^{\mu_n}\right]
    =\sum_{i=2}^{n}(-1)^i g^{\mu_1\mu_i}
    \operatorname{Tr}\left[\gamma^{\mu_2}\ldots
    \gamma^{\mu_{i-1}}\gamma^{\mu_{i+1}}\ldots\gamma^{\mu_n}\right],
\end{equation*}
где в оставшемся следе пропущены $\gamma^{\mu_1}$ и $\gamma^{\mu_i}$. Формула получается,
если переносить $\gamma^{\mu_1}$ вправо антикоммутатором
$\gamma^{\mu_1}\gamma^{\mu_i}=2g^{\mu_1\mu_i}-\gamma^{\mu_i}\gamma^{\mu_1}$
и воспользоваться цикличностью следа.

\subsection*{Окончание задачи 2 занятия 8}

При $n=2$ формула даёт
$\operatorname{Tr}[\gamma^\mu\gamma^\nu]=g^{\mu\nu}\operatorname{Tr}\mathbf 1=4g^{\mu\nu}$,
поэтому
\begin{equation*}
    \operatorname{Tr}[\hat a\hat b]
    =a_\mu b_\nu\operatorname{Tr}[\gamma^\mu\gamma^\nu]
    =4(a\cdot b).
\end{equation*}

\task[1]{Следы с \texorpdfstring{$\gamma^5$}{γ5}}
\label{task:gamma5_traces_sem9}

\condition
Вычислить $\operatorname{Tr}[\gamma^5]$, $\operatorname{Tr}[\gamma^5\hat a]$,
$\operatorname{Tr}[\gamma^5\hat a\hat b]$.

\solution
\subsection*{\texorpdfstring{$\operatorname{Tr}[\gamma^5]$}{Tr γ5}}

\begin{equation*}
    \operatorname{Tr}[\gamma^5]
    =i\operatorname{Tr}\left[\gamma^0\gamma^1\gamma^2\gamma^3\right]
    =4i\left[g^{01}g^{23}-g^{02}g^{13}+g^{03}g^{12}\right]=0,
\end{equation*}
так как все недиагональные компоненты метрики равны нулю.

\subsection*{\texorpdfstring{$\operatorname{Tr}[\gamma^5\hat a]$}{Tr γ5 a}}

\begin{equation*}
    \operatorname{Tr}[\gamma^5\hat a]
    =\operatorname{Tr}\left[\gamma^5a^\mu\gamma_\mu\right]
    =a^\mu\operatorname{Tr}\left[\gamma^5\gamma_\mu\right]=0:
\end{equation*}
под следом нечётное число (пять) матриц $\gamma^\mu$.

\subsection*{\texorpdfstring{$\operatorname{Tr}[\gamma^5\hat a\hat b]$}{Tr γ5 a b}}

\begin{align*}
    \operatorname{Tr}[\gamma^5\hat a\hat b]
    &=a_\mu b_\nu\,i\operatorname{Tr}\left[\gamma^0\gamma^1\gamma^2\gamma^3
    \gamma^\mu\gamma^\nu\right]\\
    &=a_\mu b_\nu\cdot i\left[
    -g^{0\mu}\operatorname{Tr}\left[\gamma^1\gamma^2\gamma^3\gamma^\nu\right]
    +g^{0\nu}\operatorname{Tr}\left[\gamma^1\gamma^2\gamma^3\gamma^\mu\right]
    \right]=0.
\end{align*}
Слагаемые с $g^{01}$, $g^{02}$, $g^{03}$ равны нулю, а каждый оставшийся след
четырёх матриц содержит только произведения $g^{12}$, $g^{13}$, $g^{23}$
и снова обращается в нуль.

\answer
\begin{equation*}
    \operatorname{Tr}[\gamma^5]=\operatorname{Tr}[\gamma^5\hat a]
    =\operatorname{Tr}[\gamma^5\hat a\hat b]=0.
\end{equation*}

\rtheory{Спиноры}

\subsection*{Уравнение для амплитуды}

\begin{equation*}
    \left(i\gamma^\mu\partial_\mu-m\right)\psi=0,
    \qquad
    \psi=u(\vec p)e^{-ipx},
    \qquad
    H\psi=E\psi,
\end{equation*}
\begin{equation*}
    \left(\vec\alpha\vec p+\beta m\right)u=Eu,
    \qquad
    \vec\alpha=
    \begin{pmatrix}
        0&\vec\sigma\\
        \vec\sigma&0
    \end{pmatrix},
    \qquad
    \beta=
    \begin{pmatrix}
        I&0\\
        0&-I
    \end{pmatrix}.
\end{equation*}
Записываем $u$ через два двухкомпонентных столбца $u_A$, $u_B$:
\begin{equation*}
    \left(
    \begin{pmatrix}
        0&\vec\sigma\vec p\\
        \vec\sigma\vec p&0
    \end{pmatrix}
    +
    \begin{pmatrix}
        m&0\\
        0&-m
    \end{pmatrix}
    \right)
    \begin{pmatrix}
        u_A\\ u_B
    \end{pmatrix}
    =
    \begin{pmatrix}
        m&\vec\sigma\vec p\\
        \vec\sigma\vec p&-m
    \end{pmatrix}
    \begin{pmatrix}
        u_A\\ u_B
    \end{pmatrix}
    =E
    \begin{pmatrix}
        u_A\\ u_B
    \end{pmatrix}.
\end{equation*}
Отсюда
\begin{equation*}
    \left\{
    \begin{aligned}
        \vec\sigma\vec p\,u_B&=(E-m)\,u_A,\\
        \vec\sigma\vec p\,u_A&=(E+m)\,u_B.
    \end{aligned}
    \right.
\end{equation*}

\subsection*{Решения с положительной и отрицательной энергией}

\begin{enumerate}[label=\arabic*)]
    \item $E>0$:
    \begin{equation*}
        u_A^{1,2}=\chi^{1,2},
        \qquad
        \chi^1=
        \begin{pmatrix}
            1\\0
        \end{pmatrix},
        \quad
        \chi^2=
        \begin{pmatrix}
            0\\1
        \end{pmatrix},
        \qquad
        u_B^{1,2}=\frac{\vec\sigma\vec p}{E+m}\,\chi^{1,2}.
    \end{equation*}
    \item $E<0$:
    \begin{equation*}
        u_A^{3,4}=\frac{-\vec\sigma\vec p}{|E|+m}\,\chi^{1,2},
        \qquad
        u_B^{3,4}=\chi^{1,2}.
    \end{equation*}
\end{enumerate}
Решения с отрицательной энергией описывают античастицы:
\begin{equation*}
    u^{3,4}(-\vec p)\,e^{-i(-p)x}=v^{2,1}(\vec p)\,e^{ipx}.
\end{equation*}

\subsection*{Частицы и античастицы}

\rconcept{Частицы}:
\begin{equation*}
    \underset{s=1/2}{
    \begin{pmatrix}
        \chi^1\\[0.3em]
        \dfrac{\vec\sigma\vec p}{E+m}\chi^1
    \end{pmatrix}},
    \qquad
    \underset{s=-1/2}{
    \begin{pmatrix}
        \chi^2\\[0.3em]
        \dfrac{\vec\sigma\vec p}{E+m}\chi^2
    \end{pmatrix}}.
\end{equation*}
\rconcept{Античастицы}:
\begin{equation*}
    \underset{s=1/2}{
    \begin{pmatrix}
        \dfrac{-\vec\sigma\vec p}{|E|+m}\chi^2\\[0.6em]
        \chi^2
    \end{pmatrix}},
    \qquad
    \underset{s=-1/2}{
    \begin{pmatrix}
        \dfrac{-\vec\sigma\vec p}{|E|+m}\chi^1\\[0.6em]
        \chi^1
    \end{pmatrix}}.
\end{equation*}

\remark
Здесь выписаны решения $u^{4}$ и $u^{3}$ с импульсом $\vec p$.
Спинор античастицы с физическим импульсом $\vec p$ по формуле выше
равен $v(\vec p)=u(-\vec p)$, то есть отличается \rnote{знаком перед
$\vec\sigma\vec p$}:
\begin{equation*}
    v^{1}(\vec p)=
    \begin{pmatrix}
        \dfrac{\vec\sigma\vec p}{E+m}\chi^2\\[0.6em]
        \chi^2
    \end{pmatrix}\sqrt{E+m},
    \qquad
    v^{2}(\vec p)=
    \begin{pmatrix}
        \dfrac{\vec\sigma\vec p}{E+m}\chi^1\\[0.6em]
        \chi^1
    \end{pmatrix}\sqrt{E+m}.
\end{equation*}

\subsection*{Нормировка}

Все спиноры умножаются на $\sqrt{E+m}$. При $\vec p=0$ это $\sqrt{2m}$,
и спиноры, делённые на $\sqrt{2m}$, совпадают с базисными столбцами:
\begin{align*}
    e^-&:&
    \begin{pmatrix}
        1\\0\\0\\0
    \end{pmatrix}
    &\longleftrightarrow
    \begin{pmatrix}
        \chi^1\\0
    \end{pmatrix},
    &
    \begin{pmatrix}
        0\\1\\0\\0
    \end{pmatrix}
    &\longleftrightarrow
    \begin{pmatrix}
        \chi^2\\0
    \end{pmatrix},\\[0.5em]
    e^+&:&
    \begin{pmatrix}
        0\\0\\0\\1
    \end{pmatrix}
    &\longleftrightarrow
    \begin{pmatrix}
        0\\\chi^2
    \end{pmatrix},
    &
    \begin{pmatrix}
        0\\0\\1\\0
    \end{pmatrix}
    &\longleftrightarrow
    \begin{pmatrix}
        0\\\chi^1
    \end{pmatrix}.
\end{align*}

\subsection*{Нерелятивистская частица}

Пусть $\vec p\parallel z$. Тогда
\begin{equation*}
    \frac{\vec\sigma\vec p}{E+m}
    =\frac{\sigma_zp}{E+m}
    =
    \begin{pmatrix}
        1&0\\
        0&-1
    \end{pmatrix}
    \frac{p}{E+m}.
\end{equation*}
При $v\ll c$ имеем $p\simeq mv$, $E+m\simeq2m$:
\begin{align*}
    \chi^1&:\quad
    \frac{p}{E+m}
    \begin{pmatrix}
        1\\0
    \end{pmatrix}
    \simeq\frac{mv}{2m}
    \begin{pmatrix}
        1\\0
    \end{pmatrix},\\
    \chi^2&:\quad
    \frac{p}{E+m}
    \begin{pmatrix}
        0\\-1
    \end{pmatrix}
    \simeq\frac{mv}{2m}
    \begin{pmatrix}
        0\\-1
    \end{pmatrix}.
\end{align*}
С восстановленной скоростью света $v\to v/c$.

\subsubsection*{Электрон}

\begin{equation*}
    \frac{u^{1}(\vec p)}{\sqrt{2m}}
    \simeq
    \begin{pmatrix}
        \chi^1\\[0.2em]
        \dfrac{v}{2c}\sigma_z\chi^1
    \end{pmatrix}
    =
    \begin{pmatrix}
        1\\0\\ \dfrac{v}{2c}\\0
    \end{pmatrix},
    \qquad
    \frac{u^{2}(\vec p)}{\sqrt{2m}}
    \simeq
    \begin{pmatrix}
        \chi^2\\[0.2em]
        \dfrac{v}{2c}\sigma_z\chi^2
    \end{pmatrix}
    =
    \begin{pmatrix}
        0\\1\\0\\ -\dfrac{v}{2c}
    \end{pmatrix}.
\end{equation*}

\subsubsection*{Позитрон}

Для физического импульса $\vec p\parallel +z$ берём $v(\vec p)=u(-\vec p)$,
поэтому знак перед $\sigma_z$ положительный:
\begin{equation*}
    \frac{v^{1}(\vec p)}{\sqrt{2m}}
    \simeq
    \begin{pmatrix}
        \dfrac{v}{2c}\sigma_z\chi^2\\[0.4em]
        \chi^2
    \end{pmatrix}
    =
    \begin{pmatrix}
        0\\ -\dfrac{v}{2c}\\0\\1
    \end{pmatrix},
    \qquad
    \frac{v^{2}(\vec p)}{\sqrt{2m}}
    \simeq
    \begin{pmatrix}
        \dfrac{v}{2c}\sigma_z\chi^1\\[0.4em]
        \chi^1
    \end{pmatrix}
    =
    \begin{pmatrix}
        \dfrac{v}{2c}\\0\\1\\0
    \end{pmatrix}.
\end{equation*}
Решения с отрицательной энергией $u^{4,3}(\vec p)$ отличаются от этих
знаком перед $v/(2c)$.

\rtheory{Уравнения для \texorpdfstring{$u$}{u} и \texorpdfstring{$v$}{v}}

В обозначении $\hat p=p^\mu\gamma_\mu$:
\begin{align*}
    (\hat p-m)\,u(\vec p)&=0,\\
    (-\hat p-m)\,u^{3,4}(-\vec p)&=0 \quad\text{--- ветвь с $E<0$},\\
    (\hat p+m)\,v(\vec p)&=0 \quad\text{--- для античастицы}.
\end{align*}
\rconcept{Дираковски сопряжённые} спиноры $\bar u$, $\bar v$ --- строки:
\begin{equation*}
    \bar u=u^\dagger\gamma^0,
    \qquad
    \bar v=v^\dagger\gamma^0.
\end{equation*}
Нормировка:
\begin{equation*}
    \left\{
    \begin{aligned}
        \bar u^{(s)}u^{(s)}&=2m,\\
        \bar v^{(s)}v^{(s)}&=-2m,
    \end{aligned}
    \right.
    \qquad
    s=1,2\ \text{--- спиновый индекс}.
\end{equation*}
Для $\vec p=p\,\vec e_z$
\begin{equation*}
    u^{(s)}=
    \begin{pmatrix}
        \chi^{(s)}\\[0.2em]
        \dfrac{\vec\sigma\vec p}{E+m}\chi^{(s)}
    \end{pmatrix}\sqrt{E+m}
    =
    \begin{pmatrix}
        \chi^{(s)}\\[0.2em]
        \dfrac{\sigma_zp_z}{E+m}\chi^{(s)}
    \end{pmatrix}\sqrt{E+m}.
\end{equation*}

\remark[Домашнее задание.]
\begin{enumerate}[label=\arabic*)]
    \item Вывести уравнения Дирака для $\bar u$ и $\bar v$.
    \item Найти $\bar u^{(s)}$ для $\vec p=p\,\vec e_z$ и проверить нормировку
    $\bar u^{(s)}u^{(s)}=2m$.
\end{enumerate}

\rtheory{Сводка: следы и суммы по спинам}

\remark
Раздел дополняет разобранное у доски. Здесь собраны соотношения,
которые понадобятся при вычислении $|\mathcal M|^2$ для частиц со спином.
Все формулы проверены прямым вычислением в представлении Паули--Дирака.

\subsection*{Общие свойства следа}

\begin{equation*}
    \operatorname{Tr}[A+B]=\operatorname{Tr}A+\operatorname{Tr}B,
    \qquad
    \operatorname{Tr}[AB]=\operatorname{Tr}[BA],
    \qquad
    \operatorname{Tr}[ABC]=\operatorname{Tr}[CAB]=\operatorname{Tr}[BCA].
\end{equation*}
Переставлять под следом можно только \rnote{циклически}:
в общем случае $\operatorname{Tr}[ABC]\ne\operatorname{Tr}[ACB]$.

\subsection*{След нечётного числа \texorpdfstring{$\gamma$}{γ}-матриц}

\begin{equation*}
    \operatorname{Tr}\left[\gamma^{\mu_1}\gamma^{\mu_2}\ldots\gamma^{\mu_{2n+1}}\right]=0.
\end{equation*}
Доказательство: вставим $\mathbf 1=(\gamma^5)^2$, перенесём один $\gamma^5$
циклически в конец и затем пронесём его налево через все $2n+1$ матриц,
каждый раз меняя знак:
\begin{align*}
    \operatorname{Tr}\left[\gamma^{\mu_1}\ldots\gamma^{\mu_{2n+1}}\right]
    &=\operatorname{Tr}\left[\gamma^5\gamma^5\gamma^{\mu_1}\ldots\gamma^{\mu_{2n+1}}\right]
    =\operatorname{Tr}\left[\gamma^5\gamma^{\mu_1}\ldots\gamma^{\mu_{2n+1}}\gamma^5\right]\\
    &=(-1)^{2n+1}\operatorname{Tr}\left[\gamma^5\gamma^5\gamma^{\mu_1}\ldots\gamma^{\mu_{2n+1}}\right]
    =-\operatorname{Tr}\left[\gamma^{\mu_1}\ldots\gamma^{\mu_{2n+1}}\right].
\end{align*}
След равен минус самому себе, значит, он равен нулю.
По той же причине $\operatorname{Tr}[\gamma^5\gamma^\mu]=0$
и $\operatorname{Tr}[\gamma^5\gamma^\mu\gamma^\nu\gamma^\rho]=0$.

\subsection*{Следы чётного числа \texorpdfstring{$\gamma$}{γ}-матриц}

\begin{align*}
    \operatorname{Tr}\mathbf 1&=4,\\
    \operatorname{Tr}\left[\gamma^\mu\gamma^\nu\right]&=4g^{\mu\nu},\\
    \operatorname{Tr}\left[\gamma^\mu\gamma^\nu\gamma^\rho\gamma^\sigma\right]
    &=4\left(g^{\mu\nu}g^{\rho\sigma}-g^{\mu\rho}g^{\nu\sigma}
    +g^{\mu\sigma}g^{\nu\rho}\right).
\end{align*}
В свёртке с векторами:
\begin{align*}
    \operatorname{Tr}[\hat a\hat b]&=4(a\cdot b),\\
    \operatorname{Tr}[\hat a\hat b\hat c\hat d]
    &=4\left[(a\cdot b)(c\cdot d)-(a\cdot c)(b\cdot d)+(a\cdot d)(b\cdot c)\right].
\end{align*}
След не меняется при записи матриц в обратном порядке:
\begin{equation*}
    \operatorname{Tr}\left[\gamma^{\mu_1}\gamma^{\mu_2}\ldots\gamma^{\mu_n}\right]
    =\operatorname{Tr}\left[\gamma^{\mu_n}\ldots\gamma^{\mu_2}\gamma^{\mu_1}\right].
\end{equation*}

\subsection*{Следы с \texorpdfstring{$\gamma^5$}{γ5}}

\begin{align*}
    \operatorname{Tr}\gamma^5&=0,\\
    \operatorname{Tr}\left[\gamma^5\gamma^\mu\gamma^\nu\right]&=0,\\
    \operatorname{Tr}\left[\gamma^5\gamma^\mu\gamma^\nu\gamma^\rho\gamma^\sigma\right]
    &=-4i\varepsilon^{\mu\nu\rho\sigma},
\end{align*}
где $\varepsilon^{\mu\nu\rho\sigma}$ --- полностью антисимметричный тензор,
\rassumption{$\varepsilon^{0123}=+1$}, $\varepsilon_{0123}=-1$.
При другом соглашении о знаке $\varepsilon^{0123}$ знак в последней формуле
меняется. Полезны также свёртки
\begin{equation*}
    \varepsilon^{\mu\nu\rho\sigma}\varepsilon_{\mu\nu\rho\sigma}=-24,
    \qquad
    \varepsilon^{\alpha\beta\mu\nu}\varepsilon_{\alpha\beta\rho\sigma}
    =-2\left(\delta^\mu_\rho\delta^\nu_\sigma-\delta^\mu_\sigma\delta^\nu_\rho\right).
\end{equation*}

\subsection*{Свёртки}

К результатам задачи~1 занятия~8 добавляется
случай четырёх множителей:
\begin{equation*}
    \gamma^\mu\hat a\hat b\hat c\hat d\gamma_\mu
    =2\left(\hat d\hat a\hat b\hat c+\hat c\hat b\hat a\hat d\right).
\end{equation*}

\subsection*{Эрмитово сопряжение}

Из $(\gamma^0)^\dagger=\gamma^0$, $(\gamma^k)^\dagger=-\gamma^k$ и антикоммутации
\begin{equation*}
    \gamma^0\gamma^{\mu\dagger}\gamma^0=\gamma^\mu,
    \qquad
    \gamma^0\gamma^{5\dagger}\gamma^0=-\gamma^5.
\end{equation*}
Для произвольной матрицы $\Gamma$ вводят $\bar\Gamma=\gamma^0\Gamma^\dagger\gamma^0$.
Тогда $\left(\bar u_1\Gamma u_2\right)^*=\bar u_2\bar\Gamma u_1$,
и, в частности, $\overline{\gamma^\mu}=\gamma^\mu$,
$\overline{\gamma^\mu\gamma^\nu}=\gamma^\nu\gamma^\mu$.

\subsection*{Суммы по спинам}

Для спиноров $u^{(s)}$ и $v^{(s)}$, нормированных множителем $\sqrt{E+m}$,
\begin{equation*}
    \sum_{s=1,2}u^{(s)}(p)\,\bar u^{(s)}(p)=\hat p+m,
    \qquad
    \sum_{s=1,2}v^{(s)}(p)\,\bar v^{(s)}(p)=\hat p-m.
\end{equation*}
Отсюда получается правило вычисления квадрата амплитуды:
сумма по спинам сводится к следу
\begin{equation*}
    \sum_{s,s'}\left[\bar u^{(s')}(p')\,\Gamma_1\,u^{(s)}(p)\right]
    \left[\bar u^{(s')}(p')\,\Gamma_2\,u^{(s)}(p)\right]^*
    =\operatorname{Tr}\left[\Gamma_1(\hat p+m)\bar\Gamma_2(\hat p'+m)\right].
\end{equation*}
Для античастицы соответствующий множитель $\hat p+m$ заменяется на $\hat p-m$.
При усреднении по начальным спинам неполяризованных частиц добавляется множитель
$1/2$ на каждый начальный фермион.

\end{document}
